\documentclass[a4paper,10pt]{article}
\usepackage{latexsym}
\usepackage[empty]{fullpage}
\usepackage{titlesec}
\usepackage{marvosym}
\usepackage[usenames,dvipsnames]{color}
\usepackage{verbatim}
\usepackage{enumitem}
\usepackage[colorlinks=true,urlcolor=blue]{hyperref}
\usepackage{fancyhdr}
\usepackage[margin=1cm,top=1.2cm]{geometry}

\pagestyle{fancy}
\fancyhf{}
\fancyfoot{}
\renewcommand{\headrulewidth}{0pt}
\renewcommand{\footrulewidth}{0pt}

% Sections formatting
\titleformat{\section}{\vspace{-4pt}\scshape\raggedright\large}{}{0em}{}[\color{black}\titlerule \vspace{-5pt}]

% Ensure that generate pdf is machine readable/ATS parsable
\ifdefined\pdfgentounicode
\pdfgentounicode=1
\fi

%-------------------- CV START --------------------
\begin{document}

\begin{center}
    {\Huge \scshape Jaydip Singh, Ph.D.} \\ \vspace{5pt}
  \href{mailto:jaydip.singh@gmail.com}{Jaydip.Singh@gmail.com} $|$ +91-6306311397 $|$ Mumbai, India \\
    \href{https://www.linkedin.com/in/jaydip-singh-15714546}{LinkedIn} $|$
  \href{https://github.com/JaydipSingh}{GitHub} $|$
  \href{https://github.com/JaydipSingh/portfolio}{Portfolio} $|$
  \href{https://inspirehep.net/authors/1591917}{Research Profile}
\end{center}

\section{Professional Summary}
Ph.D.-qualified Machine Learning Scientist and Lead GenAI Engineer with 10+ years of combined academic and industry experience spanning scientific computing, software/hardware engineering, and end-to-end AI/ML systems. Specializes in designing and deploying production-grade agentic AI on AWS --- using Amazon Bedrock, Bedrock AgentCore (agent runtimes, memory, multi-agent orchestration), and the Strands Agents framework --- to scale enterprise solutions across business domains. Deep expertise in multi-agent architectures, multi-LLM/SLM model selection, multi-index semantic retrieval and context engineering, RAG, knowledge graphs, and semantic/metadata layers (data catalog, lineage, ontology), with strong command of ML tradeoffs between deterministic and reasoning-driven logic, benchmarking/evaluation, and MLOps/observability. Lead Engineer (GenAI) at Apexon, delivering client-facing, enterprise-scale solutions on AWS. Also an independent researcher on efficient small and multimodal language models.

Independent research: two 2026 preprints on efficient small and multimodal language models trained on constrained hardware --- a 23M physics SLM that outperforms GPT-2 Large (774M) by 69\%, and a multimodal medical-physics visual-QA model reaching 89\% on a simulation-generated benchmark. Portfolio: \href{https://github.com/JaydipSingh/portfolio}{github.com/JaydipSingh/portfolio}.


\section{Technical Skills}
\begin{itemize}[leftmargin=0.5cm]
  \item \textbf{AI / GenAI / Agentic:} Amazon Bedrock, Bedrock AgentCore (agent runtimes, memory, multi-agent orchestration), Strands Agents framework, Azure OpenAI, LangChain, LangGraph, MCP, RAG, multi-agent systems, multi-LLM/SLM model selection, context engineering, multi-index semantic retrieval, prompt engineering, embeddings/vector retrieval, knowledge graphs, ontology/semantic layer, synthetic data generation, LLM benchmarking \& evaluation.
  \item \textbf{ML / Analytics:} Machine Learning, Deep Learning, SageMaker, forecasting (LSTM, P10/P50/P90 probabilistic), feature engineering, CNNs, Monte Carlo simulation, MLOps / model observability, statistical modeling.
  \item \textbf{Programming:} Python (NumPy, Pandas, PyTorch, TensorFlow, Scikit-learn), C/C++, SQL, Bash/Shell.
  \item \textbf{Cloud / Platform:} AWS Lambda, API Gateway, DynamoDB, S3, Cognito, Secrets Manager, WAF, CloudWatch, EventBridge, Step Functions, SNS, ECS/Fargate, ECR, CodeBuild, SSM, KMS, Redshift, SageMaker, Glue ETL, Bedrock, Bedrock AgentCore; Azure Functions, Azure AI Search.
  \item \textbf{DevOps / IaC:} AWS CDK, Terraform, Docker (ARM64), GitHub Actions, CI/CD, OIDC, IAM (least-privilege), observability (CloudWatch).
  \item \textbf{Frontend / Backend / Data:} React, TypeScript, FastAPI, Flask, REST/WebSocket APIs, Node.js; PostgreSQL, Chroma, Pinecone, Neo4j AuraDB, Redshift, S3 Data Lake.
  \item \textbf{Simulation/Design:} GEANT4 (detector development), COMSOL (chip/semiconductor simulation), SolidWorks (hardware design), LaTeX.
\end{itemize}

\section{Selected Projects}
\begin{itemize}[leftmargin=0.5cm]
  \item \textbf{SLM Physics Assistant --- Domain-Specialized Small Language Model} (Rust, Python, PyTorch, MPS)
  \begin{itemize}[leftmargin=0.45cm]
    \item Built an end-to-end, reproducible stack --- Rust BPE tokenizer, custom decoder-only transformer (RoPE, RMSNorm, SwiGLU), LoRA fine-tuning, and multi-metric evaluation --- trained entirely on a single Apple M3 Pro laptop.
    \item A 23M-parameter physics-specialized model scores 0.55 on a physics QA benchmark, outperforming GPT-2 Large (774M) by 69\% despite being 33$\times$ smaller.
    \item Diagnosed and fixed an early end-of-sequence collapse failure mode via EOS-free document packing; added automated tests and CI. Research paper (preprint; arXiv forthcoming).
    \item \href{https://github.com/JaydipSingh/slm-physics-assistant}{GitHub Repo}
  \end{itemize}
  \item \textbf{PhysVision-SLM --- Multimodal Medical-Physics Visual QA} (PyTorch, SigLIP, Qwen2, Transformers, MPS)
  \begin{itemize}[leftmargin=0.45cm]
    \item Trained a multimodal SLM entirely on simulation-generated data (an analytical image simulator provides exact ground truth, so no human annotation is required).
    \item A frozen SigLIP + Qwen2-0.5B configuration reaches 89.0\% on a 35,596-pair visual-QA benchmark (vs. 25\% random); a controlled lite-vs-full comparison attributes the gap to the vision encoder.
    \item Built the dataset generator, two-stage training, and log-probability-ranking evaluation with diagnostic checks. Research paper (preprint; arXiv forthcoming).
    \item \href{https://github.com/JaydipSingh/physvision-slm}{GitHub Repo}
  \end{itemize}
  \item \textbf{Agentic Video Generator} (FastAPI, LangGraph, Redis, Postgres, ffmpeg, Streamlit, React)
  \begin{itemize}[leftmargin=0.45cm]
    \item Built an end-to-end multi-agent video generation pipeline that transforms story prompts into rendered MP4 outputs via a checkpointed LangGraph workflow.
    \item Implemented asynchronous API + worker architecture with Redis queue and Postgres-backed job state for reliable long-running generation tasks.
    \item Added node-level checkpoint persistence (Redis cache + Postgres durable store) with resume endpoint support to recover from mid-pipeline failures.
    \item Delivered two production-ready UIs (Streamlit and React) with submission forms, status polling, resume action, and in-app video preview.
    \item \href{https://github.com/JaydipSingh/agentic-video-generator}{GitHub Repo}
  \end{itemize}
\end{itemize}

\section{Professional Experience}

\textbf{Lead Engineer (GenAI) – Apexon (Client: AWS),} Mumbai, India  \hfill Dec 2025 – Present
\begin{itemize}[leftmargin=0.5cm]
\item \textbf{Congress Operations Management Platform} (07/2026 – Present): Led end-to-end design and delivery of a cloud-based operations platform unifying attendee processing, booth staffing, report matching, and operational analytics. Built Bedrock-based agents (Supervisor on Bedrock AgentCore + Strands Agents framework), a React/TypeScript frontend and Python serverless backend, with Terraform/AWS CDK infrastructure, Cognito/IAM role-based security, and Docker/ECR + GitHub Actions deployment pipelines (ARM64).
\item \textbf{AI-Powered Clinical Trial Supply Chain Predictive Analytics Platform} (04/2026 – Present): Architected a 7-layer AWS serverless platform predicting drug demand and detecting stockout/expiry risk for a pharmaceutical client. Built SageMaker LSTM demand forecasting (P10/P50/P90, MAPE $<$20\%), 5 specialist Amazon Bedrock agents, event-driven ingestion (EventBridge/Lambda/Glue), and automated compliance (TOR) reporting via Step Functions + Bedrock Claude --- cutting reporting from days to under 60 seconds, saving 20+ hours/week and preventing \$400K+/year in losses per trial program.
\item Applied ML judgment on deterministic vs. reasoning-driven logic and multi-LLM/SLM model selection, with systematic benchmarking of agent accuracy, latency, and cost against enterprise SLAs; operationalized MLOps/observability (drift detection, retraining triggers, CloudWatch metrics).
\end{itemize}

\textbf{Gen AI Technical Lead (IC) – FutureSoft (Client: PwC BSC)}, Mumbai, India \hfill Jul 2025 - Nov 2025
\begin{itemize}[leftmargin=0.5cm]
  \item Leading Generative AI projects for PwC, including enterprise-scale chatbot/agent development, RAG-based automation, and client-facing AI solutions.
  \item Architecting and deploying AI/ML pipelines with focus on LLMs, LangChain, and cloud-based MLOps practices.
  \item Mentoring cross-functional teams on GenAI adoption, responsible AI practices, and production-grade deployments.
  \item Collaborating with stakeholders to translate business requirements into scalable AI-driven solutions.
\end{itemize}

\textbf{Postdoctoral Research Scientist (Data Scientist) – UC Davis}, USA \hfill Jan 2024 – Dec 2024
\begin{itemize}[leftmargin=0.5cm]
  \item Developed AI agents for real-time data monitoring and anomaly detection using RAG, BERT, Flan-T5, and GPT.
  \item Integrated Pinecone vector DB for log ingestion, enabling advanced reasoning and context retention.
  \item Applied ML algorithms (PCA, DBSCAN, CNNs) for image classification and 3D reconstruction.
  \item Built Monte Carlo–based models for optimizing detector performance.
\end{itemize}

\textbf{Senior Software Engineer – Nextgen Clearing Pvt Ltd.}, Mumbai, India \hfill Nov 2021 – May 2023
\begin{itemize}[leftmargin=0.5cm]
  \item Engineered telecom billing systems using Python, C++, and PL/SQL.
  \item Refactored legacy modules, boosting performance and maintainability.
  \item Coordinated development and debugging with cross-functional teams.
\end{itemize}

\textbf{AI Research Collaborator – University of Lucknow \& Fermilab}, USA \hfill Mar 2019 – Oct 2021
\begin{itemize}[leftmargin=0.5cm]
  \item Developed CNN/GenAI models for particle physics image reconstruction (LArTPC).
  \item Optimized classification and detector performance using deep learning.
  \item Performed large-scale data analysis using PyTorch and Scikit-learn.
\end{itemize}

\textbf{Research Fellow (Scientific Software Developer) – India-based Neutrino Observatory}, Lucknow University \hfill Aug 2014 – Mar 2019
\begin{itemize}[leftmargin=0.5cm]
  \item Created GEANT4-based simulation software for experimental physics.
  \item Integrated ML models to enhance detector signal recognition and efficiency.
  \item Deployed on Fermi Cloud (AWS/GCP-like) infrastructure for distributed computing.
\end{itemize}

\section{Education}
\begin{itemize}[leftmargin=0.5cm]
  \item \textbf{Ph.D. in Particle Physics}, University of Lucknow, India \hfill 2019 – 2022
  \item \textbf{M.Sc. in Physics}, Indian Institute of Technology Bombay, India \hfill 2011 – 2013
  \item \textbf{B.Sc. in Physics, Mathematics \& Computer Science}, Ewing Christian College, Allahabad \hfill 2007 – 2010
\end{itemize}


\section{Research Publications}
Extensive publication record in high-energy physics and AI/ML research. \\
Full list available at: \href{https://inspirehep.net/authors/1591917}{inspirehep.net/authors/1591917}

\subsection*{Preprints / Under Review (Independent Research)}
\begin{itemize}[leftmargin=0.5cm]
  \item \textbf{A Reproducible Systems Stack for Domain-Specialized Small Language Models: Tokenizer, Training, and Evaluation on Constrained Hardware} (2026). Preprint; arXiv forthcoming. Code: \href{https://github.com/JaydipSingh/slm-physics-assistant}{github.com/JaydipSingh/slm-physics-assistant}.
  \item \textbf{Simulation-Generated Data for Multimodal Small Language Models: A Constrained-Hardware Recipe for Medical Physics Visual Question Answering} (2026). Preprint; arXiv forthcoming. Code: \href{https://github.com/JaydipSingh/physvision-slm}{github.com/JaydipSingh/physvision-slm}.
\end{itemize}

\subsection*{Selected AI/ML-Enhanced Publications (Key Contributions)}
\begin{itemize}[leftmargin=0.5cm]
  \item \textbf{Separation of track- and shower-like energy deposits in ProtoDUNE-SP using a convolutional neural network}, Eur. Phys. J. C 82 (2022). \\
  \textit{Technologies: Convolutional Neural Network (CNN), supervised classification, image-based detector data.} \\
  Contributed to integration of CNN model into the reconstruction workflow and optimization with real and simulated data. \\
  \href{https://inspirehep.net/literature/2060793}{inspirehep.net/2060793}
  
  \item \textbf{Highly-parallelized simulation of a pixelated LArTPC on a GPU}, JINST 18 (2023). \\
  \textit{Technologies: GPU parallelization (CUDA), high-performance computing, simulation pipelines.} \\
  Developed and optimized GPU-accelerated simulation kernels and integrated with experiment frameworks. \\
  \href{https://inspirehep.net/literature/2620145}{inspirehep.net/2620145}
  
  \item \textbf{Neutrino interaction vertex reconstruction in DUNE with integrated U-ResNet neural network} (2025). \\
  \textit{Technologies: U-ResNet (deep residual U-Net), hit-level classification, hybrid ML pipeline integration.} \\
  Integrated U-ResNet into vertex reconstruction pipeline, ensuring efficient interfacing and runtime optimization. \\
  \href{https://inspirehep.net/literature/2878306}{inspirehep.net/2878306}
\end{itemize}


\section{Certifications \& Test Scores}
\begin{itemize}[leftmargin=0.5cm]
  \item AI Engineer Production Track: Deploy LLMs \& Agents at Scale (Udemy)
  \item AI Engineer Agentic Track: The Complete Agent \& MCP Course (Udemy)
  \item Generative AI with Large Language Models (AWS, DeepLearning.AI)
  \item IIT-JAM (Physics): All India Rank 263 (2011)
  \item NGPE (IAPT): Top 10 percentile (2010)
\end{itemize}

\section{Awards \& Fellowships}
\begin{itemize}[leftmargin=0.5cm]
  \item Rajendra Raja Fellowship – Fermilab, USA
  \item INO Fellowship – DST, India
\end{itemize}

\section{Extracurricular Activities}
\begin{itemize}[leftmargin=0.5cm]
  \item Winner: 4x100m Relay (IIT-Bombay \& ECC), Basketball/Volleyball tournaments
  \item Finalist: 400m Sprint, Annual college meets (2008–2010)
\end{itemize}

\end{document}
